\begin{activity}{Aliasing Demo}

\section*{Purpose}

This demonstration extends the basic ideas of Fourier transforms developed in Demo~1 to explore the impacts of sampling on what information can be recovered from series data.

The notebook progresses deliberately from:
\begin{itemize}
  \item synthetic geological models with known truth,
  \item explore impact of sampling size,
  \item see how aliasing can cause misinterpretation.
\end{itemize}

\section*{In-Class Demonstration Outline}

A suggested walkthrough is:

\begin{enumerate}
  \item \textbf{Synthetic folds:} increase station spacing until aliasing occurs.
  \item \textbf{Synthetic Basin--Range:} observe dominance of regional signal.
  \item \textbf{Spectra:} show migration of power with coarser sampling.
\end{enumerate}


\section*{Sampling and Aliasing in Space}

A spatially varying gravity field $g(x)$ can be represented as a sum of spatial
frequencies (wavenumbers):
\begin{equation}
g(x) = \int \hat{g}(k)\, e^{i 2\pi k x}\, \mathrm{d}k ,
\end{equation}
where $k$ is the spatial frequency (cycles per unit distance) and
the corresponding wavelength is $\lambda = 1/k$.

If the field is sampled at spacing $\Delta x$, the Nyquist wavelength is
\begin{equation}
\lambda_N = 2 \Delta x .
\end{equation}

Wavelengths shorter than $\lambda_N$ cannot be represented correctly.
Instead, their energy is mapped to longer wavelengths:
this is \emph{aliasing}.

\subsection*{Key property of aliasing}

Aliasing is not noise and not a numerical error.  The Fourier transform recovers a \emph{coherent but incorrect} spectrum.  Once aliasing occurs, the original short-wavelength information is irretrievably lost.

To put it simply, there are four consequences of aliasing:
\begin{itemize}
  \item You can't get what you want
  \item You get something else instead
  \item You don't know that you've done it
  \item You can't fix it after it is done
\end{itemize}

\section*{Synthetic Fold and Basin--Range Models}

To make aliasing unambiguous, the notebook first uses synthetic geological models.

\subsection*{Anticline--syncline model}

The folded interface is defined as
\begin{equation}
z(x) = z_0 + A \sin\!\left(\frac{2\pi x}{\lambda}\right),
\end{equation}
with wavelength $\lambda = 20$~km.
A pedagogical forward model is used in which the gravity response is proportional
to interface depth:
\begin{equation}
g(x) \propto z(x).
\end{equation}

This model is not physically exact, but it preserves the correct relationship between
geological wavelength and gravity wavelength.
It is sufficient for illustrating sampling and aliasing.

\subsection*{Basin--and--Range style model}

A more realistic gravity field is constructed as the superposition of
two spatial scales:
\begin{equation}
z(x) = z_0
+ A_s \sin\!\left(\frac{2\pi x}{\lambda_s}\right)
+ A_d \sin\!\left(\frac{2\pi x}{\lambda_d}\right),
\end{equation}
where:
\begin{itemize}
  \item $\lambda_s \sim 20$~km represents individual basins and ranges,
  \item $\lambda_d \sim 120$~km represents deep crustal or lithospheric structure.
\end{itemize}

The resulting gravity profile mimics real Basin-and-Range Bouguer gravity:
short-wavelength basin anomalies ride on a strong long-wavelength regional field.

\section*{Detrending and Its Limitations}

Polynomial detrending removes long-wavelength components of a signal.
In spectral terms, it acts as a crude high-pass filter.
However:
\begin{itemize}
  \item detrending is not physically motivated,
  \item high-order detrending can remove real geological signal,
  \item detrending does not correct aliasing.
\end{itemize}

The notebook uses detrending only as a stepping stone toward a physically based
regional correction.


\section*{Key Ideas}

\begin{itemize}
  \item Station spacing sets a hard limit on geological resolution.
  \item Aliasing produces coherent but incorrect structure.
  \item Isostatic correction is a physical hypothesis expressed as a Fourier filter.
  \item Spectral methods do not protect against poor sampling.
  \item Interpretation begins with survey design, not inversion.
\end{itemize}

\end{activity}